\section{Package Architecture}

The \texttt{bifurcationatlas} package provides a structured framework for constructing bifurcation atlases consisting of a central bifurcation diagram together with a collection of surrounding region illustrations. The package separates configuration, data management, rendering, layout generation, and referencing into distinct layers.

\subsection{Overview}

The workflow of the package can be summarized as

\begin{equation*}
    \text{Configuration}
\rightarrow
\text{Atlas Environment}
\rightarrow
\text{Region Registry}
\rightarrow
\text{Rendering}
\rightarrow
\text{Referencing}.
\end{equation*}
The user interacts only with the high-level commands

\begin{itemize}
\item \lstinline|\begin{bifurcationatlas}...\end{bifurcationatlas}|,
\item \lstinline|\baDeclareRegion|,
\item \lstinline|\baNewPic|,
\item \lstinline|\baAtlasFigure|,
\item \lstinline|\baSetup|.
\end{itemize}

All indexing, caption management, and reference generation are performed automatically.

\subsection{Configuration Layer}

Global and local appearance settings are managed through a PGF key family.

Global settings are specified using

\begin{lstlisting}[language=LaTeX]
\baSetup{
region scale=1.2,
caption label=roman
}
\end{lstlisting}

while local settings may be supplied directly to the atlas environment

\begin{lstlisting}[language=LaTeX]
\begin{bifurcationatlas}[caption label=alpha]
...
\end{bifurcationatlas}
\end{lstlisting}

The available configuration parameters are

\begin{center}
\begin{tabular}{ll}
\lstinline|region scale| & scale factor of region illustrations\\
\lstinline|region radius| & radial distance of regions\\
\lstinline|figure scale| & scale factor of the background figure\\
\lstinline|caption style| & font style of captions\\
\lstinline|caption shift| & vertical caption displacement\\
\lstinline|caption label| & numbering format of region captions
\end{tabular}
\end{center}

\subsection{Atlas Environment Layer}

The environment

\begin{lstlisting}[language=LaTeX]
\begin{bifurcationatlas}
...
\end{bifurcationatlas}
\end{lstlisting}

forms the central container of the package.

Internally, the environment

\begin{enumerate}
\item creates a local scope,
\item applies optional local settings,
\item increments the atlas counter,
\item resets region counters,
\item initializes region storage,
\item captures the user-defined atlas label,
\item opens the floating atlas container.
\end{enumerate}

All state information is local to the current atlas and is automatically discarded when the environment ends.

\subsection{Floating Layer}

The package defines the floating environment

\begin{lstlisting}[language=LaTeX]
bifurcationatlasfloat
\end{lstlisting}

using the \texttt{newfloat} package.

This layer is responsible only for

\begin{itemize}
\item page placement,
\item numbering of atlases,
\item generation of the list of bifurcation atlases.
\end{itemize}

All atlas-specific logic is handled separately by the wrapper environment.

\subsection{Region Registry}

Regions are declared using

\begin{lstlisting}[language=LaTeX]
\baDeclareRegion[caption]{name}{color list}
\end{lstlisting}

Each declaration stores

\begin{itemize}
\item a symbolic region identifier,
\item an optional caption,
\item a color specification.
\end{itemize}

The registry therefore acts as a database describing the regions that appear in an atlas.

\subsection{Pic System}

Reusable TikZ illustrations are introduced through

\begin{lstlisting}[language=LaTeX]
\baNewPic{name}{tikz code}
\end{lstlisting}

A pic defines the graphical representation of a region independently of any specific atlas.

Before rendering, region-specific color information is automatically expanded and injected into the pic.

This separation allows the same pic definition to be reused for many different atlases.

\subsection{Rendering Layer}

Rendering is performed by

\begin{lstlisting}[language=LaTeX]
\baRenderRegion{pic}{region}
\end{lstlisting}

which combines

\begin{itemize}
\item the selected pic,
\item the region color specification,
\item the current scale settings.
\end{itemize}

The rendering layer is independent of layout and positioning.

\subsection{Layout Layer}

The radial arrangement of regions is generated by

\begin{lstlisting}[language=LaTeX]
\baRadialAtlas
\end{lstlisting}

which

\begin{enumerate}
\item computes the number of regions,
\item distributes them uniformly on a circle,
\item creates labels and references,
\item places the rendered region graphics.
\end{enumerate}

The layout engine therefore controls only geometric placement and does not depend on the actual graphical content.

\subsection{Atlas Figure Layer}

The command

\begin{lstlisting}[language=LaTeX]
\baAtlasFigure
\end{lstlisting}

provides the highest-level user interface.

It combines

\begin{itemize}
\item a background bifurcation diagram,
\item an optional background caption,
\item the radial region layout.
\end{itemize}

This command is intended to be the primary mechanism for constructing complete bifurcation atlases.

\subsection{Referencing System}

Each atlas receives a user-defined label

\begin{lstlisting}[language=LaTeX]
\begin{bifurcationatlas}
\label{ba:example}
...
\end{bifurcationatlas}
\end{lstlisting}

The package automatically derives labels for all contained regions.

For example

\begin{lstlisting}[language=LaTeX]
ba:example.1
ba:example.2
ba:example.3
\end{lstlisting}

are generated automatically.

The displayed reference format depends on the selected caption style.

Examples include


\begin{equation*}
    1.a,\qquad
1.iv,\qquad
1.3
\end{equation*}
for alphabetic, Roman, and Arabic numbering respectively.

In addition, subreferences are stored in the auxiliary file and may be accessed through

\begin{lstlisting}[language=LaTeX]
\basubref{ba:example.3}
\end{lstlisting}

which returns only the local caption identifier.

\subsection{Design Principles}

The package follows four fundamental design principles:

\begin{enumerate}
\item Separation of configuration and rendering.
\item Reusable graphical region definitions.
\item Automatic generation of references.
\item Independence of layout and graphical content.
\end{enumerate}

This layered structure allows the package to remain extensible while maintaining a simple user-facing interface.

\subsection{PIC System, Rendering Pipeline, and Radial Layout Engine}

This section explains the core visualization pipeline of the bifurcation atlas package. It consists of four tightly connected layers:

\begin{itemize}
\item PIC system (definition of reusable graphical components),
\item rendering layer (instantiation of TikZ pictures),
\item layout engine (geometric placement),
\item atlas renderer (integration of structure, labels, and output).
\end{itemize}

% =========================================================
\subsubsection{PIC System: Reusable TikZ Components}
% =========================================================

\begin{lstlisting}[language=LaTeX]
\newcommand{\baNewPic}{[2]}{%
  \tikzset{
    #1/.pic={\baPrepareColors #2}
  }%
}
\end{lstlisting}

\paragraph{Purpose}

This command defines reusable TikZ \emph{pic objects}, which are parametrized drawing templates.

\paragraph{Syntax}

\begin{itemize}
\item \texttt{\#1}: name of the pic (identifier)
\item \texttt{\#2}: TikZ drawing code
\end{itemize}

\paragraph{Mechanism}

The command registers a new TikZ pic via:

\begin{lstlisting}[language=LaTeX]
#1/.pic={...}
\end{lstlisting}

This allows later usage as:

\begin{lstlisting}[language=LaTeX]
\pic{<name>}
\end{lstlisting}

\paragraph{Color preparation hook}

\lstinline|\baPrepareColors| is executed inside the pic definition so that:

\begin{itemize}
\item region-specific color macros are available locally,
\item each pic has access to the current atlas color state.
\end{itemize}

\subsubsection{High-Level Atlas Composition Command: \texttt{\textbackslash baAtlasFigure}}

The command

\begin{lstlisting}[language=LaTeX]
\NewDocumentCommand{\baAtlasFigure}{O{} m m m}{...}
\end{lstlisting}

is the primary user-level interface of the bifurcation atlas package. It combines three independent subsystems:

\begin{itemize}
\item a background bifurcation diagram (static image),
\item an optional caption system with labeling,
\item a radial overlay generated by \lstinline|\baRadialAtlas|.
\end{itemize}

It therefore acts as a \textbf{composition layer} that integrates data, layout, and visualization.



\paragraph{Function signature}

The command takes four arguments:

\begin{itemize}
\item \texttt{\#1} (optional): caption text for the background diagram
\item \texttt{\#2}: filename of the background image
\item \texttt{\#3}: list of regions (input for radial atlas)
\item \texttt{\#4}: TikZ pic template used for rendering regions
\end{itemize}



\paragraph{Step 1: Background diagram embedding}

\begin{lstlisting}[language=LaTeX]
\node (bg) at (0,0)
  {\includegraphics[scale=\ba@figure@scale]{#2}};
\end{lstlisting}

This places the bifurcation diagram as a TikZ node named \texttt{bg}. The scaling factor is controlled globally via:

\begin{lstlisting}[language=LaTeX]
\ba@figure@scale
\end{lstlisting}

This ensures consistent visual scaling across all atlas figures.



\paragraph{Step 2: Optional caption handling}

The optional argument \texttt{\#1} controls whether a caption is inserted. The test:

\begin{lstlisting}[language=LaTeX]
\if\relax\detokenize{#1}\relax
\end{lstlisting}

checks if the argument is empty.



\subparagraph{Case A: no caption}

If no caption is provided:

\begin{lstlisting}[language=LaTeX]
\def\ba@offset{0}
\end{lstlisting}

This disables caption-related indexing in the radial overlay.



\subparagraph{Case B: caption present}

If a caption exists:

\begin{lstlisting}[language=LaTeX]
\def\ba@offset{1}
\end{lstlisting}

and additional metadata processing is triggered:

\paragraph{Counter updates}

\begin{lstlisting}[language=LaTeX]
\refstepcounter{baSubRegion}
\refstepcounter{baRegion}
\end{lstlisting}

These ensure:

\begin{itemize}
\item a new globally referenceable region is created,
\item LaTeX referencing via \lstinline|\ref| is activated.
\end{itemize}



\paragraph{Label generation}

\begin{lstlisting}[language=LaTeX]
\label{\ba@atlaslabel.\arabic{baRegion}}
\end{lstlisting}

This assigns a hierarchical label of the form:

\begin{center}
\texttt{atlaslabel.1}, \texttt{atlaslabel.2}, \dots
\end{center}



\paragraph{Manual auxiliary file entry}

\begin{lstlisting}[language=LaTeX]
\protected@write\@auxout{}{%
  \string\newlabel{sub@\ba@atlaslabel.\arabic{baRegion}}%
  {{\thebaSubRegion}{\thepage}}%
}
\end{lstlisting}

This explicitly writes a secondary label for sub-referencing. It stores:

\begin{itemize}
\item the formatted sub-region identifier,
\item the page number.
\end{itemize}

This enables custom commands such as \lstinline|\basubref|.



\paragraph{Caption rendering}

\begin{lstlisting}[language=LaTeX]
\node[anchor=north]
  at ([yshift=-\ba@caption@shift]bg.south)
  {\ba@caption@style (\baFormatCaptionLabel)~#1};
\end{lstlisting}

This places the caption below the background image using:

\begin{itemize}
\item vertical offset control (\lstinline|\ba@caption@shift|),
\item formatting style (\lstinline|\ba@caption@style|),
\item dynamic numbering (\lstinline|\baFormatCaptionLabel|).
\end{itemize}



\paragraph{Step 3: Radial overlay integration}

\begin{lstlisting}[language=LaTeX]
\baRadialAtlas[\ba@offset]{#3}{#4}
\end{lstlisting}

This call integrates the radial atlas system on top of the background figure.

The parameter \lstinline|\ba@offset| controls whether caption-related indexing is included in the radial layout.



\paragraph{Architectural interpretation}

This command acts as a \textbf{composition orchestrator}:

\begin{itemize}
\item \textbf{Input layer}: image + region data
\item \textbf{Metadata layer}: labels, counters, auxiliary references
\item \textbf{Presentation layer}: caption + scaling
\item \textbf{Visualization layer}: radial atlas overlay
\end{itemize}

It is therefore the highest-level abstraction in the package, binding together all lower-level rendering and indexing mechanisms into a single user-facing interface.

% =========================================================
\subsubsection{Rendering Layer: Instantiating a Region}
% =========================================================

\begin{lstlisting}[language=LaTeX]
\newcommand{\baRenderRegion}[2]{%
\scalebox{\ba@region@scale}{%
\begin{tikzpicture}
  \pic[ba/region/#2]{#1};
\end{tikzpicture}}%
}
\end{lstlisting}

\paragraph{Purpose}

This command renders a single region using a predefined pic template.

\paragraph{Parameters}

\begin{itemize}
\item \texttt{\#1}: pic name (drawing template)
\item \texttt{\#2}: region identifier (used for styling)
\end{itemize}

\paragraph{Pipeline}

\begin{enumerate}
\item A TikZ picture is created.
\item A pic is inserted with style \texttt{ba/region/<region>}.
\item The result is scaled via \lstinline|\scalebox|.
\end{enumerate}

\paragraph{Key idea}

This separates:

\begin{itemize}
\item \textbf{geometry} (pic definition),
\item \textbf{data} (region ID),
\item \textbf{presentation scaling}.
\end{itemize}



% =========================================================
\subsubsection{Layout Engine: Radial Positioning}
% =========================================================

\begin{lstlisting}[language=LaTeX]
\newcommand{\baPlaceRadialNode}[3]{%
  \pgfmathsetmacro{\ba@angle}{90-(#1-1)*360/#2}%
  \node at (\ba@angle:\ba@region@radius) {#3};
}
\end{lstlisting}

\paragraph{Purpose}

This command computes angular positions for radial layout placement.

\paragraph{Parameters}

\begin{itemize}
\item \texttt{\#1}: index of current region
\item \texttt{\#2}: total number of regions
\item \texttt{\#3}: content to place
\end{itemize}

\paragraph{Angle computation}

The formula:

\begin{lstlisting}[language=LaTeX]
90 - (i-1)*360/n
\end{lstlisting}

ensures:

\begin{itemize}
\item first element starts at top (90°),
\item elements are evenly distributed,
\item layout is circular.
\end{itemize}

\paragraph{Geometric placement}

The syntax:

\begin{lstlisting}[language=LaTeX]
(\ba@angle:\ba@region@radius)
\end{lstlisting}

is polar coordinate notation in TikZ.



% =========================================================
\subsubsection{Radial Atlas Renderer: Full Integration}
% =========================================================

\begin{lstlisting}[language=LaTeX]
\NewDocumentCommand{\baRadialAtlas}{O{0} m m}{%
\baRequireAtlas
...
}
\end{lstlisting}

\paragraph{Purpose}

This is the main orchestration layer that combines:

\begin{itemize}
\item data iteration,
\item counter management,
\item label generation,
\item rendering,
\item caption formatting.
\end{itemize}



\paragraph{Step 1: compute number of regions}

\begin{lstlisting}[language=LaTeX]
\foreach \r [count=\i] in {#2}{%
  \xdef\baLastIndex{\i}%
}
\end{lstlisting}

This loop determines the total number of regions for angular spacing.



\paragraph{Step 2: region iteration}

Each region is processed in a second loop:

\begin{lstlisting}[language=LaTeX]
\foreach \r [count=\i] in {#2}{%
\end{lstlisting}



\paragraph{Step 3: counter and reference management}

\begin{lstlisting}[language=LaTeX]
\refstepcounter{baSubRegion}
\refstepcounter{baRegion}
\label{\ba@atlaslabel.\arabic{baRegion}}
\end{lstlisting}

This ensures:

\begin{itemize}
\item each region is numbered,
\item labels are referenceable,
\item hyperref integration works automatically.
\end{itemize}

Additionally:

\begin{lstlisting}[language=LaTeX]
\protected@write\@auxout{}{...}
\end{lstlisting}

manually writes a secondary label (\texttt{sub@...}) for sub-referencing.



\paragraph{Step 4: caption retrieval}

\begin{lstlisting}[language=LaTeX]
\edef\baTmpCaption{\baGetRegionCaption{\r}}
\end{lstlisting}

This stores region-specific metadata for later conditional display.



\paragraph{Step 5: radial placement and rendering}

\begin{lstlisting}[language=LaTeX]
\baPlaceRadialNode{\i}{\baLastIndex}{...}
\end{lstlisting}

Each region is placed on the circle and rendered using:

\begin{itemize}
\item \lstinline|\baRenderRegion| for the TikZ pic,
\item optional caption display,
\item tabular layout for vertical stacking.
\end{itemize}



\subsubsection{Architectural Summary}

This subsystem implements a full rendering pipeline:

\begin{itemize}
\item \textbf{PIC layer}: reusable graphical primitives
\item \textbf{render layer}: instantiation of region visuals
\item \textbf{layout layer}: geometric positioning in polar coordinates
\item \textbf{atlas layer}: integration of structure, labels, and output
\end{itemize}

The design separates concerns cleanly:

\begin{itemize}
\item data (regions),
\item geometry (pics),
\item layout (radial mapping),
\item metadata (labels and captions).
\end{itemize}